\section{Norma algebraica ternaria y radios óptimos de recuperación}
\label{lib02:norma-ternaria}

El radio de recuperación de las dos memorias contiene el factor
\(\sqrt{73}\). Su origen puede seguirse sin introducir una magnitud
objetivo: procede del Gram del lector nonádico y, en otra realización del
ciclo ternario ya construido, es la escala de una semejanza. Su cuadrado,
\(73\), es el índice de la inclusión reticular de rango dos. Estas
apariciones se reúnen mediante operaciones explícitas;
no se identifican por una semejanza decimal con alfa. La comparación
posterior con \(10^4\alpha_A\) se tratará como un problema inverso,
con sus datos, sus transformaciones permitidas y sus conclusiones exactas.

La materia prima conserva su producción APP--TRIT--TPK y el corte
posterior descrito en \cref{lib01:registro}. El registro entero, la
orientación de los desplazamientos, la carga y los lectores se reciben
antes de calcular distancias. El caso canónico utiliza \(8+1\) y el
denominador nueve. El parámetro \(a>0\) que se introduce a continuación
es una colección algebraica posterior, no una nueva semilla ni un parámetro
físico seleccionado por alfa. Los errores pertenecen al canal de lectura;
no se los identifica con eventos o historias admisibles.

\subsection{Colección de dos memorias y su Gram}
\label{lib02:familia-memoria}

En \(V=\mathbb R^{12}\), sea \((Sx)_i=x_{i+1}\), con índices módulo
doce. Para \(a>0\) y \(j=3,4\), definimos
\begin{equation}
 T_{j,a}=\frac{aI+S^j}{a+1},\qquad
 D_{j,a}=\frac{\sqrt a(S^j-I)}{a+1},\qquad
 M_ax=(D_{3,a}x,D_{4,a}T_{3,a}x).
 \label{lib02:lectores}
\end{equation}
La segunda memoria se produce después de la primera compresión. Archivar
ambas no equivale a reinyectar el primer complemento en una evolución
reversible de dos canales. Sean
\[
 p(a)=a^2+a+1,\quad A(a)=a(a^2+1),\quad B(a)=a^2,
 \qquad\Pi=I-\frac{\mathbf1\mathbf1^*}{12},\quad P_u=I-\Pi.
\]

\begin{proposition}[Balance y forma de diferencias]
\label{lib02:gram}
Para todo \(a>0\),
\[
 T_{j,a}^*T_{j,a}+D_{j,a}^*D_{j,a}=I,\qquad
 G_a:=M_a^*M_a=I-(T_{4,a}T_{3,a})^*(T_{4,a}T_{3,a}),
\]
y
\begin{align}
 (a+1)^4G_a={}&4ap(a)I
 -A(a)(S^3+S^4+S^8+S^9)\notag\\
 &-B(a)(S+S^5+S^7+S^{11}).
 \label{lib02:gram-explicito}
\end{align}
Equivalentemente,
\begin{align}
 (a+1)^4\|M_ah\|^2={}&A(a)\sum_i
 \bigl((h_i-h_{i+3})^2+(h_i-h_{i+4})^2\bigr)\notag\\
 &+B(a)\sum_i
 \bigl((h_i-h_{i+1})^2+(h_i-h_{i+5})^2\bigr).
 \label{lib02:energia-aristas}
\end{align}
\end{proposition}
\begin{proof}
La ortogonalidad de \(S\) da
\[
 T_{j,a}^*T_{j,a}
 =\frac{(a^2+1)I+a(S^j+S^{-j})}{(a+1)^2},\qquad
 D_{j,a}^*D_{j,a}
 =\frac{a(2I-S^j-S^{-j})}{(a+1)^2}.
\]
Su suma es \(I\); aplicar la identidad a las dos etapas cancela el
término intermedio \(T_{3,a}^*T_{3,a}\). Como todos los operadores
son polinomios de \(S\), el producto de los dos Grams individuales
tiene diagonal \((a^2+1)^2\), términos simples de peso
\(a(a^2+1)\) y cruzados de peso \(a^2\). La identidad
\((a+1)^4-(a^2+1)^2=4a(a^2+a+1)\) da
\eqref{lib02:gram-explicito}. Expandir las diferencias cuadradas da
\eqref{lib02:energia-aristas}; cada arista no orientada de cada paso
aparece una sola vez y el grado ponderado es \(4A+4B=4ap(a)\).
\end{proof}

En \(a=8\) se obtiene
\begin{align}
 6561G_8={}&2336I-520(S^3+S^4+S^8+S^9)\notag\\
 &-64(S+S^5+S^7+S^{11}),
 \label{lib02:gram-nonadico}
\end{align}
cuya primera fila es
\((2336,-64,0,-520,-520,-64,0,-64,-520,-520,0,-64)/6561\).
También
\[
 T_{4,8}T_{3,8}=\frac{64I+8S^3+8S^4+S^7}{81}.
\]
Los cuatro desplazamientos distintos producen la contribución diagonal
\[
 1-\frac{64^2+8^2+8^2+1}{81^2}
 =\frac{2336}{6561}
 =\frac{32(8^2+8+1)}{9^4}.
\]

\begin{theorem}[Núcleo, espectro e inversa completa]
\label{lib02:inversa-memorias}
Para todo \(a>0\), \(\ker M_a=\mathbb R\mathbf1\). Los valores
propios de \(G_a\), indexados por los modos duodécimos \(m\), son
\[
\begin{array}{c|c|c}
 m&\lambda_m(G_a)&\text{multiplicidad}\\\hline
 0&0&1\\
 3,9&2a/(a+1)^2&2\\
 4,8&3a/(a+1)^2&2\\
 6&4a/(a+1)^2&1\\
 1,5,7,11&a(5a^2+4a+5)/(a+1)^4&4\\
 2,10&a(7a^2+2a+7)/(a+1)^4&2
\end{array}
\]
En particular,
\begin{equation}
 L_a^{\rm rec}=(G_a+P_u)^{-1}M_a^*,\qquad
 L_a^{\rm rec}M_a=\Pi,\qquad
 \|L_a^{\rm rec}\|=\frac{a+1}{\sqrt{2a}}.
 \label{lib02:inversa}
\end{equation}
\end{theorem}
\begin{proof}
Si \(M_ah=0\), la primera memoria implica \(S^3h=h\), luego
\(T_{3,a}h=h\); la segunda implica \(S^4h=h\). Como
\(S=S^4(S^3)^{-1}\), el vector es uniforme. La recíproca es
inmediata. En el modo \(m\), el balance da
\[
 \lambda_m(G_a)=1-
 \frac{[a^2+1+2a\cos(\pi m/2)]
 [a^2+1+2a\cos(2\pi m/3)]}{(a+1)^4}.
\]
Los doce pares de cosenos dan la tabla. Respecto de
\(2a/(a+1)^2\), las diferencias de los dos últimos valores son
\[
 \frac{3a(a^2+1)}{(a+1)^4}>0,\qquad
 \frac{a(5a^2-2a+5)}{(a+1)^4}>0.
\]
Los otros valores positivos son múltiplos mayores. Por tanto el mínimo
en \(\mathbf1^\perp\) es \(2a/(a+1)^2\). El operador
\(G_a+P_u\) vale uno sobre la uniforme y es positivo invertible en
su complemento. La fórmula escrita es la inversa de mínimos cuadrados
y su norma es el inverso de la menor singularidad no nula de \(M_a\).
\end{proof}

En \(a=1\), \(T_{3,1}=(I+S^3)/2\) es singular, pero
\eqref{lib02:inversa} sigue siendo válida. La singularidad de una
compresión individual no elimina la recuperación mediante sus dos
memorias. En \(a=8\), la norma de la inversa es \(9/4\).

\subsection{Distancias mínimas: todos los registros enteros}
\label{lib02:redes}

Los dominios algebraicos de decodificación son
\[
 \mathcal C_0(a)=M_a\mathbb Z^{12},\qquad
 \Lambda_q=\{k\in\mathbb Z^{12}:\mathbf1^*k=q\},\qquad
 \mathcal C_q(a)=M_a\Lambda_q.
\]
La primera imagen se identifica con la red \(\Pi\mathbb Z^{12}\)
o con clases módulo \(\mathbb Z\mathbf1\); la segunda conserva una
carga entera fijada. Son dominios que contienen las representaciones enteras
del tipo indicado. Su uso no afirma que cada candidato corresponda a una
historia HMT producida.

\begin{lemma}[Corte mínimo del grafo de memoria]
\label{lib02:cortes}
Todo corte propio no vacío de las doce posiciones tiene al menos cuatro
aristas de pasos \(1,5\) y al menos cuatro de pasos \(3,4\). Su peso
en \eqref{lib02:energia-aristas} es, por tanto, al menos
\(4A(a)+4B(a)=4ap(a)\); un corte de una posición alcanza esa cota.
\end{lemma}
\begin{proof}
Cada paso \(1\) y \(5\) forma un ciclo conexo de doce posiciones;
cada ciclo cruza un corte no trivial al menos dos veces. Para los pasos
\(3,4\), si el conjunto no es invariante por ninguno, cada colección
aporta al menos dos aristas. Si es invariante por \(S^3\), es unión
de clases módulo tres; cada uno de los cuatro triángulos de \(S^4\)
lo cruza dos veces, dando ocho aristas. Si es invariante por \(S^4\),
es unión de clases módulo cuatro; cada uno de los tres ciclos de longitud
cuatro de \(S^3\) lo cruza al menos dos veces, dando al menos seis.
La invariancia simultánea obligaría a invariancia por \(S\), imposible
para un conjunto propio no vacío. Todos los pesos son positivos; el
corte de una sola posición tiene exactamente cuatro aristas de cada
peso.
\end{proof}

\begin{theorem}[Distancias exactas y minimizadores de carga fija]
\label{lib02:distancias}
Para todo \(a>0\) y toda carga entera \(q\),
\begin{equation}
 d_0(a)^2=\min_{h\notin\mathbb Z\mathbf1}\|M_ah\|^2
 =\frac{4ap(a)}{(a+1)^4},\qquad
 d_q(a)^2=\frac{8ap(a)}{(a+1)^4}.
 \label{lib02:distancias-exactas}
\end{equation}
En el primer mínimo \(h\) recorre \(\mathbb Z^{12}\); el segundo
es la distancia mínima entre palabras distintas de \(\mathcal C_q(a)\).
Las diferencias que alcanzan el segundo mínimo son exactamente
\(e_i-e_j\) con \(j-i\equiv2,6,10\pmod{12}\).
\end{theorem}
\begin{proof}
Para un entero \(h\) no constante,
\((h_i-h_j)^2\geq|h_i-h_j|\). Cada diferencia absoluta se
descompone como suma de los cortes de los conjuntos de nivel
\(\{i:h_i>t\}\). Hay al menos un nivel no trivial; por
\cref{lib02:cortes}, la energía ponderada es al menos \(4ap(a)\).
El vector \(e_i\) alcanza esa energía. Dividir por \((a+1)^4\)
prueba el primer mínimo sobre todos los enteros, no sólo sobre
indicadores de subconjuntos.

Para carga fija, la diferencia \(h\ne0\) es entera y de suma cero.
Su norma cuadrada es par porque \(h_i^2\equiv h_i\pmod2\). Si
\(\|h\|^2\geq4\), el mínimo espectral da
\[
 \|M_ah\|^2\geq\frac{8a}{(a+1)^2}
 >\frac{8ap(a)}{(a+1)^4},
\]
ya que \((a+1)^2-p(a)=a>0\). Sólo queda
\(\|h\|^2=2\), es decir, \(h=e_i-e_j\). Su energía es
\(2((G_a)_{ii}-(G_a)_{ij})\). Los elementos no diagonales son cero
en separaciones \(2,6,10\), negativos de peso \(B(a)/(a+1)^4\)
en \(1,5,7,11\), y negativos de peso \(A(a)/(a+1)^4\) en
\(3,4,8,9\). La diagonal es \(4ap(a)/(a+1)^4\); esto prueba el
mínimo y la clasificación exhaustiva para toda la semirrecta \(a>0\).
\end{proof}

\begin{theorem}[Radios estrictos de corrección de las memorias]
\label{lib02:radios-memoria}
Para un dato \(y=M_ak+e\), con error en la norma conjunta de las dos
memorias, los radios uniformes óptimos son
\begin{equation}
 r_0(a)=\frac{\sqrt{ap(a)}}{(a+1)^2},\qquad
 r_q(a)=\frac{\sqrt{2ap(a)}}{(a+1)^2}.
 \label{lib02:radios}
\end{equation}
Si \(\|e\|<r_0(a)\), el vecino más próximo de
\(\mathcal C_0(a)\) recupera exactamente \(\Pi k\). Con carga
\(q\) conocida y \(\|e\|<r_q(a)\), el vecino más próximo de
\(\mathcal C_q(a)\) recupera exactamente \(k\).
\end{theorem}
\begin{proof}
La red \(\Pi\mathbb Z^{12}\subset(1/12)\mathbb Z^{12}\) es
discreta y \(M_a\) está acotado inferiormente en
\(\mathbf1^\perp\). Sus imágenes son discretas y admiten vecino
más próximo. Si \(c'\) es otra palabra de código y \(d\) es la
distancia mínima correspondiente,
\(\|y-c'\|\geq d-\|e\|>\|e\|\) cuando \(\|e\|<d/2\).
El mínimo es único y la preimagen es única en cada dominio indicado.
El punto medio de dos palabras a distancia mínima está a \(d/2\)
de ambas: ningún decodificador puede distinguir uniformemente ambas
entradas desde ese mismo dato. La optimalidad se refiere a los dominios
algebraicos completos, no a un subconjunto adicional de historias que
pudiera tener una separación mayor.
\end{proof}

En el caso canónico,
\begin{equation}
 r_0(8)=\frac{2\sqrt{146}}{81}\simeq0.298346814162829,
 \qquad
 r_q(8)=\frac{4\sqrt{73}}{81}\simeq0.421926110879878.
 \label{lib02:radio73}
\end{equation}
El \(4\) procede de \(\sqrt{2\cdot8}\) al tomar la mitad de la
distancia entre registros de carga fija; \(73=8^2+8+1\) procede del
Gram; \(81=9^2\) procede de las dos normalizaciones nonádicas. El mismo
producto \(9\times9\) cuenta las celdas APP, pero esos dos papeles
no se identifican sin su aplicación tipada. Tampoco el factor cuatro
prueba por sí solo una afirmación sobre direcciones espaciales.

Si las memorias se almacenan como \(z=M_ak/\sqrt a\), la norma de
ruido cambia. En \(a=8\), los radios de ese dato racional son
\(r_0/\sqrt8=\sqrt{73}/81\) y
\(r_q/\sqrt8=\sqrt{146}/81\). Usar los radios de
\eqref{lib02:radio73} sin esta conversión es cambiar el problema.
La aplicación a \(Q_K\), \(\eta_K\), \(B_K\) y la pantalla
incidencial se ha probado en \cref{lib01:correccion-incidencia}; en
particular, el umbral exacto \(K_4=729\) no impide su restitución
discreta.

\subsection{Decodificador finito y dualidad del parámetro}
\label{lib02:computacion}

\begin{proposition}[Suelo y techo sin tabla objetivo]
\label{lib02:decodificador}
Con carga conocida, bajo el radio \(r_q(a)\), basta examinar a lo
sumo \(2^{12}=4096\) candidatos enteros. Sin carga, bajo
\(r_0(a)\), basta examinar a lo sumo \(12\cdot2^{12}=49152\)
candidatos repartidos entre doce residuos de carga. El mínimo residual
devuelve respectivamente \(k\) o su clase uniforme.
\end{proposition}
\begin{proof}
Con carga \(q\), forme
\(\widetilde k=L_a^{\rm rec}y+(q/12)\mathbf1\). Entonces
\[
 \|\widetilde k-k\|
 <\|L_a^{\rm rec}\|r_q(a)
 =\frac{\sqrt{p(a)}}{a+1}<1.
\]
Cada coordenada entera verdadera pertenece a
\(\{\lfloor\widetilde k_i\rfloor,
\lceil\widetilde k_i\rceil\}\). Se enumeran esos productos,
se eliminan los de suma distinta de \(q\) y se minimiza
\(\|M_ak'-y\|\). El candidato correcto está presente y
\cref{lib02:radios-memoria} asegura su unicidad. Coordenadas ya enteras
reducen el número de candidatos.

Sin carga, escriba \(\rho\in\{0,\ldots,11\}\) para el residuo
de la carga. Cada clase módulo \(\mathbb Z\mathbf1\) tiene un
único representante cuya suma es \(\rho\), obtenido restando
\(\lfloor q/12\rfloor\mathbf1\). Para cada \(\rho\), forme
\(\widetilde k_\rho=L_a^{\rm rec}y+(\rho/12)\mathbf1\).
En el residuo correcto,
\[
 \|\widetilde k_\rho-k_\rho\|
 <\|L_a^{\rm rec}\|r_0(a)
 =\frac{\sqrt{p(a)}}{\sqrt2(a+1)}<1.
\]
Se aplica suelo y techo, se conserva la suma \(\rho\) y se minimiza
el residual entre los doce grupos. La unicidad de la clase de código
da una única salida centrada. Estos procedimientos usan el lector y la
integralidad, no reciben \(K\) como tabla de comparación.
\end{proof}

Para \(a\) racional se puede conservar aritmética racional escribiendo
\(R_a=M_a/\sqrt a\), \(z=y/\sqrt a\). La reconstrucción centrada
es
\[
 (R_a^*R_a+P_u)^{-1}R_a^*z=L_a^{\rm rec}y.
\]
El residual puede compararse mediante sus cuadrados racionales. Estas
garantías se refieren a doce coordenadas: no afirman complejidad
polinómica para una dimensión arbitraria, autenticación criptográfica ni
ganancia cuántica. El ejemplo perturbado de
\cref{lib01:correccion-incidencia} ilustra el algoritmo canónico.

\begin{proposition}[Dualidad recíproca y máximo del radio]
\label{lib02:dualidad-a}
La sustitución \(a\mapsto1/a\) conserva el Gram, las distancias,
los radios y la norma de la inversa. Se tienen
\[
 D_{j,1/a}=D_{j,a},\qquad
 T_{j,1/a}=S^jT_{j,a}^*.
\]
Los radios máximos de esta colección se alcanzan únicamente en \(a=1\):
\[
 r_q(a)\leq\frac{\sqrt6}{4},\qquad
 r_0(a)\leq\frac{\sqrt3}{4}.
\]
\end{proposition}
\begin{proof}
Las dos identidades operatorias resultan al multiplicar numerador y
denominador por \(a\). Los Grams individuales y la fórmula
\eqref{lib02:gram} son invariantes bajo la reciprocidad. Ponga
\(t=a/(a+1)^2\in(0,1/4]\). Entonces
\[
 r_q(a)^2=2t(1-t),\qquad r_0(a)^2=t(1-t).
\]
Ambas expresiones crecen para \(0<t\leq1/4\), y \(t=1/4\)
equivale a \(a=1\). La reciprocidad deja \(t\) fijo. Al tender
\(a\) a cero o a infinito, \(t\) y los radios tienden a cero.
\end{proof}

La igualdad de Grams no significa que \(M_a=M_{1/a}\): el segundo
complemento incluye una compresión diferente. Los extremos \(a=0\)
y \(a=\infty\) no pertenecen al dominio del teorema. Optimizar la
colección en \(a=1\), o utilizar el representante recíproco \(1/8\),
no cambia retroactivamente el origen canónico \(a=8\).

\subsection{La norma en el módulo de aumento ternario}
\label{lib02:modulo-ternario}

El ciclo \(r\mapsto4r\pmod9\), con órbitas \((1,4,7)\) y
\((2,8,5)\), proporciona el módulo de aumento cuya coordenatización y forma son
\[
 C=\begin{pmatrix}0&-1\\1&-1\end{pmatrix},\qquad
 G_{A_2}=\begin{pmatrix}2&-1\\-1&2\end{pmatrix},\qquad
 C^2+C+I=0,\qquad C^TG_{A_2}C=G_{A_2}.
\]
La acción y su orientación contragrediente pertenecen a las fibras
ternarias antecedentes. Sobre ese módulo ya construido se componen
\[
 Q_-(a)=aI-C,\qquad Q_+(a)=(a+1)I+C.
\]

\begin{proposition}[Semejanza ternaria e índice entero]
\label{lib02:norma-A2}
Para \(a>0\),
\[
 Q_-Q_+=p(a)I,\qquad Q_-^{\dagger_{G_{A_2}}}=Q_+,
 \qquad Q_-^TG_{A_2}Q_-=p(a)G_{A_2},\qquad \det Q_-=p(a).
\]
En \(a=8\), ambas coordenatizaciones conjugadas son
\[
 Q_-=
 \begin{pmatrix}8&1\\-1&9\end{pmatrix},\qquad
 Q_+=\begin{pmatrix}9&-1\\1&8\end{pmatrix},\qquad Q_-Q_+=73I.
\]
Cada cociente entero es cíclico de orden 73. En el cociente de \(Q_-\),
\(C\) actúa por \(8\); en el de \(Q_+\), por
\(-9\equiv64\equiv8^{-1}\pmod{73}\).
\end{proposition}
\begin{proof}
Expandir el producto y sustituir \(C^2=-C-I\) da \(p(a)I\).
La invariancia de la forma implica
\(C^{\dagger_{G_{A_2}}}=C^{-1}=C^2=-I-C\), de donde se obtiene
el adjunto. Componer producto y adjunto da la semejanza. El determinante
de \(\bigl(\begin{smallmatrix}a&1\\-1&a+1\end{smallmatrix}\bigr)\)
es \(p(a)\). En \(a=8\), el máximo común divisor de las entradas
de cada matriz es uno y su determinante es 73; la forma de Smith es
\(\operatorname{diag}(1,73)\). La relación \(8I-C=0\) en el
primer cociente y \(9I+C=0\) en el segundo da las acciones indicadas.
Como \(8^3\equiv1\pmod{73}\) y \(8\not\equiv1\), ambas
acciones tienen orden tres y orientaciones inversas.
\end{proof}

Así, \(\sqrt{73}\) es una escala de semejanza en la métrica ternaria
heredada, no sólo una cifra dentro de un radio. Su reconocimiento complejo
posterior es \(|8-\omega|^2=73\), con
\(\omega^2+\omega+1=0\). La unidad compleja describe la realización
del operador; no fue utilizada para seleccionar el registro. El operador
\(Q_-\) no se identifica por estas igualdades con \(K\), con toda
la dinámica TPK ni con el conúcleo del lector de escala de acción.

\subsection{El ciclo completo y el origen del término fijo}
\label{lib02:ciclo-completo}

En las doce posiciones, sea \(U=S^4\). Ahora hay cuatro ciclos de
longitud tres. Sus proyectores son
\[
 P_0=\frac{I+U+U^2}{3},\qquad P=I-P_0.
\]
El rango de \(P_0\) es cuatro y el de \(P\) ocho. En particular,
\(P_0\ne P_u\), pues este último sólo conserva la uniforme global;
tampoco se confunde \(P\) con el proyector excepcional \(P_3\).
Utilizaremos ahora \(Q_{-,a}=aI-U\) y \(Q_{+,a}=(a+1)I+U\)
para los operadores del ciclo completo.

\begin{theorem}[Gram ternario con sector fijo]
\label{lib02:gram-ternario}
Se tienen
\begin{equation}
 Q_{-,a}^*Q_{-,a}
 =(a-1)^2P_0+p(a)P=p(a)I-3aP_0.
 \label{lib02:gram-Qmenos}
\end{equation}
Para \(v\ne0\), sean
\[
 w(v)=\frac{\|P_0v\|^2}{\|v\|^2},\qquad
 s_{-,a}(v)=\frac{\|Q_{-,a}v\|^2}{\|v\|^2}.
\]
Entonces \(s_{-,a}(v)=p(a)-3aw(v)\). En el caso canónico,
\begin{equation}
 Q_-^*Q_-=49P_0+73P=73I-24P_0,
 \qquad s(v)=73-24w(v)\in[49,73].
 \label{lib02:73-24w}
\end{equation}
\end{theorem}
\begin{proof}
De \(U^3=I\) y \(U^*=U^2\) resulta que \(P_0\) es el
proyector ortogonal sobre \(\ker(U-I)\). Al expandir el Gram,
\[
 Q_{-,a}^*Q_{-,a}=(a^2+1)I-a(U+U^2).
\]
La relación \(U+U^2=3P_0-I\) da
\eqref{lib02:gram-Qmenos}. Evaluar sobre \(v\) y dividir por su
norma cuadrada demuestra las demás identidades.
\end{proof}

El término \(24\) es exactamente \(3\cdot8\): es la diferencia
\(p(8)-(8-1)^2\) entre la lectura de aumento y la fija. Ya aparece
en un solo ciclo de tres posiciones. Para \(m\) ciclos, el rango fijo
es \(m\), la dimensión es \(3m\) y
\[
 \operatorname{tr}(Q_{-,a}^*Q_{-,a})=3m(a^2+1).
\]
En \(m=4,a=8\), la traza es 780 y la corrección total de la traza
es \(24\cdot4=96\). La igualdad numérica \(24=2\cdot12\) no es
el origen operatorio del coeficiente.

Las dos coordenatizaciones completas distinguen también sus determinantes. En un
solo ciclo,
\[
 \det(8I-U)=8^3-1=511=7\cdot73,
 \qquad
 \det(9I+U)=9^3+1=730=10\cdot73.
\]
El factor fijo es, respectivamente, 7 y 10; el módulo de aumento aporta
73. Sobre las doce posiciones,
\(\det(8I-U)=511^4=68184176641\). Se conserva la identidad
\begin{equation}
 73=8^2+8+1=9^2-9+1=\frac{9^3+1}{9+1}
 =\frac{729+1}{10}.
 \label{lib02:completacion73}
\end{equation}
Su interpretación como índice y factor de determinante precede a
cualquier lectura de completación o refinamiento posterior.

\subsection{Seis estados recibidos y conservación de la lectura completa}
\label{lib02:seis-estados}

Fijamos de nuevo \(a=8\) y omitimos ese subíndice en los lectores.
Los seis objetos que se comparan son
\[
 K,\quad k=\Pi K,\quad m_0=D_3K,\quad m_1=D_4T_3K,
 \quad u=P_3k,\quad q_K=Q_KK=k-u.
\]
Ninguno se selecciona a partir del objetivo del ensayo. Para dar las
evaluaciones sin sustituirlas por decimales, escribamos
\(a_v=\|v\|^2\), \(b_v=\|P_0v\|^2\). Se obtiene
\begin{equation}
\begin{array}{c|c|c}
 v&a_v&b_v\\\hline
 K&4251257&10684363/3\\
 k&11789915/12&1170761/4\\
 m_0&15413392/81&16838032/243\\
 m_1&1105759232/6561&0\\
 u&1327717/4+(568896/5)\sqrt5
   &3029657/20+(145120/3)\sqrt5\\
 q_K&1951691/3-(568896/5)\sqrt5
   &2292404/15-(145120/3)\sqrt5
\end{array}
 \label{lib02:tabla-normas}
\end{equation}
La tabla se obtiene aplicando \(P_0=(I+S^4+S^8)/3\), las fórmulas
de \(D_3,D_4T_3\) y el proyector finito de carácter
\eqref{lib01:P3} a \eqref{lib01:K}. Todas sus entradas pertenecen a
\(\mathbb Q(\sqrt5)\); las primeras cuatro son racionales. En
particular, los tres primeros pesos y lecturas son
\begin{align*}
 w(K)&=\frac{10684363}{12753771},&
 s(K)&=\frac{224866857}{4251257},\\
 w(k)&=\frac{3512283}{11789915},&
 s(k)&=\frac{776369003}{11789915},\\
 w(m_0)&=\frac{1052377}{2890011},&
 s(m_0)&=\frac{61904585}{963337}.
\end{align*}
Para \(u\) y \(q_K\), las expresiones exactas son
\(w=b_v/a_v\), \(s=73-24b_v/a_v\), con los numeradores y
denominadores de \eqref{lib02:tabla-normas}. Las evaluaciones posteriores
son
\[
\begin{array}{c|c|c}
 v&w(v)\text{, aproximadamente}&s(v)\text{, aproximadamente}\\\hline
 K&0.8377414805&52.89420446705527\\
 k&0.2979057101&65.85026295779062\\
 m_0&0.3641429046&64.26057028848679\\
 m_1&0&73\text{ exacto}\\
 u&0.4428244530591254&62.37221312658099\\
 q_K&0.1127385160275117&70.29427561533972
\end{array}
\]
La igualdad de la segunda memoria es forzada:
\(P_0D_4T_3=\sqrt8P_0(U-I)T_3/9=0\), mientras
\(\|m_1\|^2>0\). No es una coincidencia de evaluación.

\begin{proposition}[La memoria completa no modifica el peso fijo]
\label{lib02:memoria-ternaria}
La aplicación
\[
 V_2v=(t(v),m_0(v),m_1(v))
 =(T_4T_3v,D_3v,D_4T_3v)
\]
es isométrica y entrelaza \(P_0\) y \(Q_-=8I-U\) con sus
copias diagonales. Conserva por separado \(\|v\|^2\),
\(\|P_0v\|^2\) y \(\|Q_-v\|^2\). Por tanto las lecturas
conjuntas \(w\) y \(s\) no cambian.
\end{proposition}
\begin{proof}
El balance de dos etapas da \(V_2^*V_2=I\). Cada canal es un
polinomio en \(S\), por lo que conmuta con \(U,P_0,Q_-\).
Aplicar la isometría a \(v\), \(P_0v\) y \(Q_-v\) da las tres
igualdades de normas. Al dividir, la lectura total es la media de las
lecturas de los canales ponderada por sus normas cuadradas, no su media
aritmética sin pesos. Los canales nulos no aportan al cociente total.
\end{proof}

El terminal confirma materialmente esa partición. Para \(K\),
\[
 \|t(K)\|^2=\frac{25538253193}{6561},\qquad
 \|P_0t(K)\|^2=\frac{848595371}{243},
\]
\[
 w(t(K))=\frac{22912075017}{25538253193},\qquad
 s(t(K))=\frac{1314402682681}{25538253193}
 \simeq51.46799480558349.
\]
La suma de las tres normas cuadradas es \(4251257\), y la de las
normas fijas es
\((848595371+16838032)/243=10684363/3\). Para \(k=\Pi K\),
\[
 \|t(k)\|^2=\frac{16367568169}{26244},\qquad
 \|P_0t(k)\|^2=\frac{217142795}{972}.
\]
Los complementos no cambian al centrar y sus sumas recuperan las normas
de la segunda fila de \eqref{lib02:tabla-normas}. Descartar el terminal
o un complemento altera la lectura parcial; conservarlos no introduce
una corrección oculta.

En una memoria infinita completa se mantiene la misma conclusión cuando
los canales conmutan con \(U\): también el límite positivo y su raíz
conmutan con él. Sin esa conmutación se usa el transporte
\(\mathcal VP_0\mathcal V^*\) en la imagen, conforme a
\cref{lib01:isometria}, no una copia diagonal supuesta.

\subsection{Inversión, coordenatizaciones y balance bilateral}
\label{lib02:bidireccional}

\begin{proposition}[Qué conserva una inversión o un cambio de marco]
\label{lib02:inversion-carta}
Sustituir \(U\) por \(U^{-1}\) intercambia cada \(Q_{\pm,a}\)
con su adjunto y conserva \(P_0\), sus Grams, valores singulares,
determinantes y distancias inducidas. Para todo \(v\ne0\), conserva
\(w(v)\) y \(s_{-,a}(v)\). No es lo mismo que invertir
\(Q_{-,a}\), cuya inversa, para \(a\ne1\), es
\begin{equation}
 Q_{-,a}^{-1}
 =\frac1{a-1}P_0+\frac1{p(a)}Q_{+,a}P.
 \label{lib02:inversa-Q}
\end{equation}
Para una coordenatización invertible \(R\), transportando
\(U'=RUR^{-1}\), \(v'=Rv\) y
\(g'=(R^{-1})^*R^{-1}\), se conservan esas razones de normas en
la métrica \(g'\).
\end{proposition}
\begin{proof}
Como \(U^*=U^{-1}=U^2\), la inversión produce los adjuntos.
Los operadores son polinomios normales en \(U\); sus Grams dependen
de \(U+U^2\), que no cambia. La inversa escrita se verifica por
sectores: \(Q_{-,a}\) actúa por \(a-1\) en el fijo y el producto
con \(Q_{+,a}\) vale \(p(a)I\) en el aumento. Para la coordenatización,
\(Q'v'=RQv\) y \(R^*g'R=I\); numerador y denominador recuperan
exactamente sus valores originales. El adjunto respecto de \(g'\)
es también el transporte del adjunto original.
\end{proof}

La reversión del orden de las doce posiciones satisface
\(RSR^{-1}=S^{-1}\) y conserva \(P_0\). Las traslaciones \(S^j\),
incluida la antipodal \(S^6\), conmutan con \(U\); una negación
global tampoco altera las normas. Ninguna de esas operaciones convierte
por sí sola 73 en una lectura distinta.

Con la métrica euclídea fija, un marco ortogonal aplicado sólo al estado
conserva \(s_{-,a}(v)\) para todos los \(v\) si y sólo si
\(R^*P_0R=P_0\), por \eqref{lib02:gram-Qmenos} y polarización.
El grupo de estas transformaciones es
\(O(\operatorname{im}P_0)\times O(\operatorname{im}P)
\cong O(4)\times O(8)\), mayor que el conmutante de \(U\).
En cambio, un filtro de estado
\(c_0P_0+c_1P\), cuando su salida es no nula, produce
\begin{equation}
 w'=
 \frac{|c_0|^2w}{|c_0|^2w+|c_1|^2(1-w)}.
 \label{lib02:filtro-pesos}
\end{equation}
Puede cambiar la lectura. Elegir esos coeficientes desde un objetivo es
un ajuste declarado del filtro, no un cambio de coordenatización de la misma
lectura.

\begin{theorem}[Balance bilateral y recuperación explícita]
\label{lib02:balance-bilateral}
Para todo \(a>0\), en el ciclo completo se cumplen
\begin{align}
 Q_{+,a}-Q_{-,a}^*&=3P_0,&
 Q_{+,a}Q_{-,a}&=p(a)I-3P_0,\notag\\
 Q_{+,a}^*Q_{+,a}&=p(a)I+3(a+1)P_0,
 \label{lib02:dos-cartas}
\end{align}
y
\begin{equation}
 (a+1)Q_{-,a}^*Q_{-,a}
 +aQ_{+,a}^*Q_{+,a}=(2a+1)p(a)I.
 \label{lib02:balance-pesado}
\end{equation}
La aplicación
\begin{equation}
 W_av=\frac{(\sqrt{a+1}\,Q_{-,a}v,
                    \sqrt a\,Q_{+,a}v)}{\sqrt{(2a+1)p(a)}}
 \label{lib02:W}
\end{equation}
es isométrica y se invierte en su imagen mediante \(W_a^*\), de
norma uno. Para los canales sin normalizar \(x=Q_{-,a}v\),
\(y=Q_{+,a}v\),
\begin{equation}
 v=\frac{x+y}{2a+1},\qquad
 Uv=\frac{ay-(a+1)x}{2a+1}.
 \label{lib02:reconstruccion-bilateral}
\end{equation}
Un par \((x,y)\) pertenece a esa imagen sin normalizar si y sólo si
\(Q_{+,a}x=Q_{-,a}y\).
\end{theorem}
\begin{proof}
En \(Q_+-Q_-^*\) aparece \(I+U+U^2=3P_0\). Expandir
\(Q_+Q_-\), \(Q_+^*Q_+\) y usar
\(U^2=3P_0-I-U\) da \eqref{lib02:dos-cartas}. La combinación
ponderada cancela los términos de \(P_0\) en
\eqref{lib02:gram-Qmenos} y \eqref{lib02:dos-cartas}; esto prueba
\(W_a^*W_a=I\). Como \(Q_-+Q_+=(2a+1)I\), la suma de los
canales da \(v\); su combinación \(ay-(a+1)x\) da
\((2a+1)Uv\).

La condición de imagen es necesaria porque ambos operadores conmutan.
Si se cumple, defina \(v=(x+y)/(2a+1)\). Entonces
\[
 Q_-v=\frac{Q_-x+Q_-y}{2a+1}
 =\frac{(Q_-+Q_+)x}{2a+1}=x,
\]
y análogamente \(Q_+v=y\). La fórmula de la inversa adjunta es la
de toda isometría; para el par crudo, la inversa lineal
\((x,y)\mapsto(x+y)/(2a+1)\) tiene norma
\(\sqrt2/(2a+1)\) en el espacio producto.
\end{proof}

En \(a=8\), los sectores fijo y de aumento tienen lecturas
\((49,100)\) y \((73,73)\), respectivamente, y
\[
 9\|Q_-v\|^2+8\|Q_+v\|^2=17\cdot73\|v\|^2=1241\|v\|^2,
 \qquad v=\frac{Q_-v+Q_+v}{17}.
\]
Para \(s_{+,a}(v)=\|Q_{+,a}v\|^2/\|v\|^2\),
\[
 s_{+,a}=p(a)+3(a+1)w,
 \qquad (a+1)s_{-,a}+a s_{+,a}=(2a+1)p(a).
\]
La cancelación vale aun con componente fija. La composición bilateral es
una formalización sobre operadores antecedentes, no la afirmación de que
todo retorno TPK sea automáticamente este par.

Si \(\mathcal V:H\to\mathcal M\) es una memoria isométrica
completa, \(U^{\mathcal V}=\mathcal VU\mathcal V^*\) satisface
\((U^{\mathcal V})^3=\Pi_{\mathcal M}\), donde
\(\Pi_{\mathcal M}=\mathcal V\mathcal V^*\) es la identidad en
la imagen. Allí el Gram transportado es
\[
 (Q_{-,a}^{\mathcal V})^*Q_{-,a}^{\mathcal V}
 =p(a)\Pi_{\mathcal M}-3a\mathcal VP_0\mathcal V^*.
\]
Conserva las lecturas y la recuperación bilateral. Usar copias diagonales
de \(U\) sólo es legítimo cuando los canales entrelazan esa acción;
en caso contrario el transporte por \(\mathcal V\) es la definición
correcta sobre la imagen.

\subsection{El dato de alfa y las colecciones de iteraciones enteras}
\label{lib02:ensayo-inverso}

El lector completo de alfa, con acarreo, proporciona el prefijo por
bloques
\[
 (007,297,352,569,283,800,997,285,105,472,380,662,999,\ldots).
\]
La ventana de frontera cero que termina en 663 es otro objeto; no se la
intercambia con la salida completa. Para el ensayo basta el intervalo
cerrado heredado
\begin{equation}
 \frac{7297352569283800997285105472380662}{10^{36}}
 \leq\alpha_A\leq
 \frac{7297352569283800997285105472380663}{10^{36}}.
 \label{lib02:intervalo-alpha}
\end{equation}
No se identifica la constante con ninguno de sus extremos racionales.
Pongamos \(s_\alpha=10^4\alpha_A\), de modo que
\begin{equation}
 \begin{gathered}
 \ell\leq s_\alpha\leq h,\\
 \ell:=72.97352569283800997285105472380662,\\
 h:=72.97352569283800997285105472380663.
 \end{gathered}
 \label{lib02:intervalo-s}
\end{equation}
La condición inversa del lector fijo es exactamente
\begin{equation}
 s(v)=s_\alpha
 \quad\Longleftrightarrow\quad
 w(v)=\frac{73-s_\alpha}{24}
 \simeq0.00110309613175.
 \label{lib02:criterio-alpha}
\end{equation}

Ninguno de los seis estados de \cref{lib02:seis-estados} satisface esa
condición. La segunda memoria tiene \(s=73\), mientras que las otras
cinco tienen \(s<72\). Esta última desigualdad se comprueba exactamente
como \(24b_v>a_v\). En las tres filas racionales basta multiplicar
enteros; para \(u\) y \(q_K\) las diferencias son
\[
 24b_u-a_u=\frac{66073183}{20}+\frac{5235904}{5}\sqrt5>0,
\]
\[
 24b_{q_K}-a_{q_K}
 =\frac{45259241}{15}-\frac{5235904}{5}\sqrt5>0.
\]
Para la segunda, \(\sqrt5<5/2\) da ya una cota inferior positiva.
La exclusión concierne a este lector y estos estados; no declara una
ausencia global de relaciones con alfa.

Como \(0<s_\alpha<73\), el radio
\[
 r_\alpha=\frac{4\sqrt{s_\alpha}}{81}<\frac{4\sqrt{73}}{81}
\]
es una garantía conservadora válida, pero no es el radio óptimo del
lector canónico. Cualquier número entre cero y 73 daría de igual modo
una garantía menor. Si se impusiera \(a^2+a+1=s_\alpha\) cambiando
\(a\), cambiarían también \(\sqrt{2a}\) y \((a+1)^2\); no se
puede sustituir sólo el 73 y conservar 4 y 81 como si el operador no
hubiera cambiado.

\begin{theorem}[Iteración exacta del peso y recurrencia de la lectura]
\label{lib02:iteraciones}
Sea \(v_n=Q_{-,a}^nv_0\), \(a>0\), y
\(b=(a-1)^2\). Cuando \(v_n\ne0\),
\begin{equation}
 w_n=\frac{b^nw_0}{b^nw_0+p(a)^n(1-w_0)}.
 \label{lib02:pesos-iterados}
\end{equation}
Para \(a\ne1\) y \(0<w_0<1\),
\[
 \frac{w_n}{1-w_n}
 =\left(\frac b{p(a)}\right)^n\frac{w_0}{1-w_0},
 \quad w_n\downarrow0,\quad s_n\uparrow p(a),
\]
y
\begin{equation}
 s_{n+1}=p(a)+b-\frac{bp(a)}{s_n},\qquad
 s_{n+1}-s_n=\frac{(s_n-b)(p(a)-s_n)}{s_n}>0.
 \label{lib02:riccati}
\end{equation}
\end{theorem}
\begin{proof}
El operador y sus adjuntos preservan los dos sectores ortogonales. Por
\eqref{lib02:gram-Qmenos}, la norma cuadrada fija se multiplica por
\(b\) en cada paso y la complementaria por \(p(a)\). Esto da
\eqref{lib02:pesos-iterados} y la razón de pesos. Como
\(p(a)-b=3a>0\), la monotonía y el límite son estrictos en el caso
mixto. Finalmente,
\(s_n=\|v_{n+1}\|^2/\|v_n\|^2\); la sucesión de normas es
una combinación de \(b^n\) y \(p(a)^n\), que satisface
\(a_{n+2}=(p(a)+b)a_{n+1}-bp(a)a_n\). Dividir por
\(a_{n+1}\) da la recurrencia y factorizar su diferencia da
\eqref{lib02:riccati}.
\end{proof}

Si \(a=1\), el sector fijo se anula en el primer paso. Cuando el
complemento no es nulo, la lectura posterior vale exactamente tres;
un estado puramente fijo se transforma en cero y su cociente deja de
estar definido. Para \(a\ne1\), la iteración inversa satisface
\(s_{n-1}=bp(a)/(p(a)+b-s_n)\) y disminuye la lectura mixta.
Invertir la orientación \(U\mapsto U^{-1}\), en cambio, deja
invariada la sucesión de normas de las iteraciones hacia delante.
Las potencias de \(Q_{+,a}\) multiplican la razón de pesos por
\(((a+2)^2/p(a))^n>1\) y disminuyen \(s_{-,a}\) en estados
mixtos. Son filtros reales del estado, no cambios de coordenatización ni una
trayectoria temporal TPK identificada por definición.

\begin{proposition}[Ensayo completo de todas las iteraciones enteras]
\label{lib02:cruces-enteros}
Para \(a=8\), ninguna iteración entera \(Q_-^nv\), \(n\geq0\),
con \(v=K,k,m_0\), satisface \(s=s_\alpha\). Los únicos cruces
del intervalo objetivo se producen entre los índices \(21,22\),
\(14,15\) y \(15,16\), respectivamente. Para \(v=m_1\), la
lectura permanece exactamente 73 en todas las iteraciones.
\end{proposition}
\begin{proof}
Definamos los extremos racionales del peso objetivo
\[
 w_L=\frac{73-h}{24}
 =\frac{2647430716199002714894527619337}
 {2400000000000000000000000000000000},
\]
\[
 w_H=\frac{73-\ell}{24}
 =\frac{441238452699833785815754603223}
 {400000000000000000000000000000000}.
\]
Para \(w_0=A/(A+B)\), la fórmula exacta es
\(w_n=A49^n/(A49^n+B73^n)\). Los datos y los cruces son
\[
\begin{array}{c|r|r|c}
 v&A&B&N\\\hline
 K&10684363&2069408&21\\
 k&3512283&8277632&14\\
 m_0&1052377&1837634&15
\end{array}
\]
y satisfacen, en cada fila,
\[
 \frac{A49^N}{A49^N+B73^N}>w_H>w_L>
 \frac{A49^{N+1}}{A49^{N+1}+B73^{N+1}}.
\]
Son desigualdades de enteros: por ejemplo, si
\(w_H=p_H/q_H\), la primera equivale a
\((q_H-p_H)A49^N>p_HB73^N\), y la última se verifica
análogamente con \(w_L\). La sustitución de las tres filas y de
los enteros escritos prueba los cruces estrictos sin un redondeo de
alfa. La monotonía estricta de \cref{lib02:iteraciones} excluye
todos los índices anteriores y posteriores; entre \(N\) y
\(N+1\) no hay otro entero. Para \(m_1\), \(w_0=0\) permanece
cero por \eqref{lib02:pesos-iterados}.
\end{proof}

Como referencia numérica, las lecturas a ambos lados de los cruces son
\[
\begin{array}{c|c|c}
 v&s_N&s_{N+1}\\\hline
 K&72.97136264842803&72.98077012438719\\
 k&72.96167997982852&72.97426483341773\\
 m_0&72.96527892867391&72.97668298511348
\end{array}
\]
La prueba cubre todos los \(n\geq0\), no sólo una muestra finita.
La orientación inversa da los mismos cruces; las potencias positivas de
\(Q_+\) y las potencias inversas de \(Q_-\) disminuyen las
lecturas iniciales, ya menores que 72, y tampoco alcanzan el objetivo
en estas tres colecciones. Los productos mixtos \(Q_-^nQ_+^m\)
constituyen otra clase de dos índices, con razón de pesos
\((49/73)^n(100/73)^m w_0/(1-w_0)\); no quedan excluidos por
la prueba uniparamétrica anterior ni se eligen aquí índices para ajustar
alfa.

\subsection{Reciprocidad angular y corrección regional existente}
\label{lib02:correccion-regional}

Los cambios de coordenatización de \cref{lib02:inversion-carta} incluyen la
reciprocidad y la deformación elíptica ya construidas. En un plano
realizado por \(M_\eta\), la forma transportada es
\(M_\eta^{-T}M_\eta^{-1}\); la semejanza no cambia los
autovalores del ciclo. Mantener el producto euclídeo anterior después de
cambiar la coordenatización definiría otro observable. La involución
\(q_+\leftrightarrow q_-\) conserva \(A\), revierte \(C^*\)
y envía \(\eta_{\rm el}\) a \(-\eta_{\rm el}\). La relación
polar de la elipse
\(\tan\phi=e^{-\eta_{\rm el}}\tan\theta\) cambia la
parametrización de la curva, no la fase espectral de un transporte
semejante. No añade una fase pequeña a \(S^4\).

La conjugación incidencial tiene otro tipo. Su ley de bloques conserva
\[
 z_m^\dagger=z_{\rho(m)}
 +100\mathbf1_{\{c_{\rho(m)}\ne0\}}-20c_{\rho(m)},
 \qquad c_m=[C_m]_{10},
\]
con los términos de frontera antes de la reducción decimal. No es en
general una isometría del vector de bloques. Aplicarla al registro
requiere su empalme con el registro de incidencias y sus extremos; no se la
reemplaza por una corrección angular inventada.

\begin{proposition}[Corrección regional del contraángulo]
\label{lib02:retorno-regional}
Las coordenadas antecedentes de acción y ángulo cumplen
\[
 A_{\rm deg}=1000\alpha_A,\qquad
 C^*_{\rm ret,deg}=2(\eta_{\rm ret}+\alpha_A),\qquad
 C^*_{\rm pre,deg}=2(H_5+\alpha_A),
\]
\[
 \eta_{\rm ret}=H_5-\frac{90}{\pi}\mathcal C_\pi(\alpha_A),
 \quad
 \mathcal C_\pi(\alpha_A)=169\alpha_A^6+
 \frac{D_A\alpha_A^7}{1-D_A\alpha_A},\qquad
 D_A=\exp(-100\pi\alpha_A/9).
\]
La diferencia en radianes es exactamente
\[
 \theta_{C,\rm pre}-\theta_{C,\rm ret}=\mathcal C_\pi(\alpha_A).
\]
\end{proposition}
\begin{proof}
La diferencia en grados es
\(2(H_5-\eta_{\rm ret})=180\mathcal C_\pi/\pi\).
Multiplicar una sola vez por \(\pi/180\) da el resultado.
\end{proof}

Aquí \(\pi\) y \(\alpha_A\) son las coordenadas ya generadas;
\(H_5\) y los restantes lectores conservan sus definiciones
antecedentes, no una recalibración. Las evaluaciones posteriores son
\[
\begin{array}{c|c}
 \text{coordenada}&\text{radianes, aproximadamente}\\\hline
 \theta_A&0.127362829013\\
 \theta_C&0.0370662264658\\
 \theta_C/6\text{ en cada pareja antipodal}&0.00617770441097\\
 \mathcal C_\pi&2.55207421805\cdot10^{-11}
\end{array}
\]
Los canales \(q_\pm\) son atenuaciones reales positivas; su exponente
no se convierte automáticamente en una fase unitaria mediante un factor
\(i\). No se introducen divisores adicionales para aproximar las
escalas del ensayo que sigue.

\subsection{Ensayo inverso angular en las doce posiciones}
\label{lib02:ensayo-angular}

Para un objetivo general \(s\in[49,81]\), la orientación principal
del sector puro se obtiene de
\[
 |8-e^{i\theta}|^2=65-16\cos\theta,
 \qquad \theta(s)=\arccos\frac{65-s}{16}\in[0,\pi].
\]
En particular, el dato posterior \(s_\alpha\) determina
\[
 \theta_\alpha=\theta(s_\alpha),\qquad
 \delta_\alpha=\frac{2\pi}{3}-\theta_\alpha.
\]
Esta fórmula resuelve un problema inverso con una salida ya producida;
no pretende producir alfa de nuevo.

\begin{lemma}[Estructura angular del complemento ternario]
\label{lib02:J}
El operador \(J=(U-U^2)/\sqrt3\) satisface
\[
 J^*=-J,\quad J^2=-P,\quad JP_0=P_0J=0,\quad JP=PJ=J,
 \qquad U=P_0-\frac12P+\frac{\sqrt3}{2}J.
\]
\end{lemma}
\begin{proof}
El adjunto intercambia \(U,U^2\). Además,
\((U-U^2)^2=U+U^2-2I=-3P\), y
\((I+U+U^2)(U-U^2)=0\). Las demás relaciones se deducen de
\(P=I-P_0\); sumar \((U+U^2)/2\) y \((U-U^2)/2\)
da la última expresión.
\end{proof}

\begin{theorem}[Colección angular y propiedades conservadas]
\label{lib02:Utheta}
Para \(\theta\in\mathbb R\), defina
\[
 U_\theta=P_0+\cos\theta\,P+\sin\theta\,J,\qquad
 Q_{-,\theta}=8I-U_\theta.
\]
Se tienen
\[
 U_\theta^*U_\theta=I,\qquad
 U_\theta U_\varphi=U_{\theta+\varphi},\qquad U_{2\pi/3}=U,
\]
\begin{equation}
 Q_{-,\theta}^*Q_{-,\theta}
 =49P_0+(65-16\cos\theta)P.
 \label{lib02:gram-angular}
\end{equation}
Los lectores
\(T_\theta=(8I+U_\theta)/9\),
\(D_\theta=\sqrt8(U_\theta-I)/9\) conservan
\(T_\theta^*T_\theta+D_\theta^*D_\theta=I\). En cambio,
\[
 U_\theta^3=I\ \Longleftrightarrow\ 3\theta\in2\pi\mathbb Z,
 \qquad
 \det Q_{-,\theta}=7^4(65-16\cos\theta)^4.
\]
\end{theorem}
\begin{proof}
Las relaciones de \cref{lib02:J} anulan los términos entre sectores.
En el complemento \(J^2=-I\), de modo que
\(U_\theta^*U_\theta=P_0+(\cos^2\theta+\sin^2\theta)P=I\).
Multiplicar dos expresiones y utilizar las fórmulas de adición da la
ley de composición. La expresión de \(U\) en el lema da el valor
ternario. Puesto que
\(U_\theta+U_\theta^*=2P_0+2\cos\theta P\), expandir el
Gram da \eqref{lib02:gram-angular}. Las expansiones
\[
 81T_\theta^*T_\theta=65I+8(U_\theta+U_\theta^*),\qquad
 81D_\theta^*D_\theta=16I-8(U_\theta+U_\theta^*)
\]
prueban el balance. La ley de composición da
\(U_\theta^3=P_0+\cos(3\theta)P+\sin(3\theta)J\), que es
la identidad exactamente bajo la condición indicada.

Para el determinante, \(Q_{-,\theta}\) actúa por 7 en las cuatro
dimensiones fijas. En el complemento, un vector unitario \(e\) y
\(Je\) forman un par ortonormal; su complemento ortogonal es estable
bajo \(J\). Repitiendo se obtienen cuatro planos, en cada uno de los
cuales el determinante es
\((8-\cos\theta)^2+\sin^2\theta=65-16\cos\theta\).
Multiplicar las contribuciones demuestra la fórmula.
\end{proof}

\begin{corollary}[Solución inversa y separación de correcciones]
\label{lib02:fase-alpha}
En \(\operatorname{im}P\), el lector
\(Q_{-,\theta_\alpha}\) tiene razón cuadrática exactamente
\(s_\alpha\). Las orientaciones \(\pm\theta_\alpha\) dan la
misma razón y
\[
 73-s_\alpha
 =8(1-\cos\delta_\alpha+\sqrt3\sin\delta_\alpha).
\]
Se cumple \(U_{\theta_\alpha}^3\ne I\), y
\(\delta_\alpha\ne\mathcal C_\pi(\alpha_A)\).
\end{corollary}
\begin{proof}
La definición de la fase resuelve el coeficiente complementario de
\eqref{lib02:gram-angular}; el coseno es par. Sustituir
\(\theta_\alpha=2\pi/3-\delta_\alpha\) da la identidad estable
para la diferencia. Como \(72<s_\alpha<73\),
\(\pi/2<\theta_\alpha<2\pi/3\) y
\(0<3\delta_\alpha<\pi/2\); por tanto \(3\theta_\alpha\)
no es múltiplo de \(2\pi\).

La separación de correcciones puede probarse sin confiar en cifras
aproximadas. Del intervalo resulta \(73-s_\alpha>0.026\).
Si \(\delta_\alpha\leq10^{-3}\), las desigualdades
\(1-\cos t\leq t^2/2\), \(\sin t\leq t\) y
\(\sqrt3<2\) darían
\[
 73-s_\alpha\leq8(10^{-6}/2+2\cdot10^{-3})=0.016004,
\]
contradicción. Así, \(\delta_\alpha>10^{-3}\). Por otro lado,
\(0<\alpha_A<10^{-2}\), \(0<D_A<1\), implican
\[
 0<\mathcal C_\pi(\alpha_A)
 <169\cdot10^{-12}+\frac{10^{-14}}{1-10^{-2}}
 <170\cdot10^{-12}<10^{-7}.
\]
Las dos cantidades son distintas.
\end{proof}

La evaluación diagnóstica da
\[
 \delta_\alpha\simeq0.00190956706826\ {\rm rad}
 \simeq0.109410133709^\circ,\qquad
 \theta_\alpha\simeq119.890589866291^\circ.
\]
La corrección regional es aproximadamente \(7.48\cdot10^7\) veces
menor. La colección inversa conserva ortogonalidad y balance, pero pierde
el orden tres en esa fase y cambia el determinante original. Su coordenatización
real no aporta por sí sola una matriz entera ni el cociente reticular de
índice 73. El hecho de resolver el dato prescrito no demuestra
\(73=10^4\alpha_A\).

\begin{proposition}[Alcanzabilidad al conservar el estado]
\label{lib02:alcanzabilidad}
Para \(v\ne0\), con \(w=w(v)\) fijo,
\[
 s_\theta(v)=49w+(65-16\cos\theta)(1-w).
\]
Si \(w<1\), su imagen es exactamente
\([49,81-32w]\). Por tanto existe una fase que alcanza
\(s_\alpha\) si y sólo si
\[
 w\leq\frac{81-s_\alpha}{32}\simeq0.250827322099.
\]
Cuando existe, su orientación principal es
\[
 \theta_v=\arccos\left[
 \frac1{16}\left(65-\frac{s_\alpha-49w}{1-w}\right)\right].
\]
Si \(w=1\), la lectura es siempre 49 y no alcanza el objetivo.
\end{proposition}
\begin{proof}
Evaluar \eqref{lib02:gram-angular} y dividir por \(\|v\|^2\)
da la fórmula. El coseno recorre \([-1,1]\), por lo que el
coeficiente complementario recorre \([49,81]\). Su combinación con
el peso fijo da el intervalo completo, sin huecos. Resolver la igualdad
y comprobar el dominio del arcocoseno da la condición y la fase.
El caso \(w=1\) se evalúa sin dividir por \(1-w\).
\end{proof}

\begin{corollary}[Cuatro imposibilidades y dos fases posibles]
\label{lib02:seis-angulares}
Manteniendo el estado y el sector fijo, la colección angular no alcanza
\(s_\alpha\) sobre \(K,k,m_0,u\). Sí lo alcanza sobre
\(m_1\) y \(q_K\).
\end{corollary}
\begin{proof}
De la tabla exacta \eqref{lib02:tabla-normas} se obtiene
\[
\begin{array}{c|c}
 v&32b_v-9a_v\\\hline
 K&227115677/3\\
 k&2094607/4\\
 m_0&122655440/243\\
 u&37201759/20+(7859008/15)\sqrt5
\end{array}
\]
Todas las entradas son positivas. Así \(w>9/32\), y el máximo
\(81-32w\) es menor que \(72<s_\alpha\). Para \(m_1\),
\(b_{m_1}=0\) y \(a_{m_1}>0\), luego el intervalo es
\([49,81]\). Para \(q_K\),
\[
 a_{q_K}-4b_{q_K}=\frac{588839+1195712\sqrt5}{15}>0.
\]
Por tanto \(w(q_K)<1/4\) y su máximo es mayor que
\(73>s_\alpha\). En ambos casos el objetivo pertenece al
intervalo alcanzable.
\end{proof}

Las evaluaciones de los máximos y, cuando existen, las fases son
\[
\begin{array}{c|c|c}
 v&\max_\theta s_\theta(v)&\theta_v\\\hline
 K&54.1922726227&\text{no existe}\\
 k&71.4670172771&\text{no existe}\\
 m_0&69.3474270513&\text{no existe}\\
 u&66.8296175021&\text{no existe}\\
 m_1&81&119.8905898663^\circ\\
 q_K&77.3923674871&133.5296853010^\circ
\end{array}
\]
La fase del sector puro no sirve indiscriminadamente para un estado
mixto: su lectura sería
\(49w+s_\alpha(1-w)=s_\alpha-(s_\alpha-49)w\).
Los lectores originales de \(V_2\) son polinomios en \(S\), como
\(P_0,P,J,U_\theta\). Por ello entrelazan también la colección
angular, y \(s_\theta(V_2v)=s_\theta(v)\). Archivar los
complementos conserva la lectura conjunta, pero no elimina las
restricciones de alcanzabilidad.

\subsection{Resultado de la composición y alcance de los ensayos}
\label{lib02:alcance}

El factor \(73\) ha quedado determinado por tres operaciones
compatibles y diferenciadas: el Gram de dos lectores nonádicos, la
semejanza del módulo de aumento ternario y los determinantes enteros de
las dos coordenatizaciones conjugadas. El balance bilateral conserva el estado
completo y admite inversión explícita; la memoria completa conserva las
lecturas transportadas, mientras la selección de un canal puede cambiar
sus pesos. Los radios de corrección son uniformes en todos los registros
enteros del dominio algebraico y en toda la colección \(a>0\), con
pruebas de cortes y paridad que no dependen de un muestreo.

El ensayo inverso añade condiciones comprobables, no una calibración
oculta. La inversión de orientación y la covarianza métrica no cambian
la lectura. Las iteraciones enteras especificadas quedan resueltas para
todos sus índices por un cruce exacto y monotonía. La variación angular
produce una colección que realiza el objetivo en los estados permitidos,
pero el dato \(s_\alpha\) selecciona su fase y se pierde el índice
ternario original. Su identificación con una realización previa de
\(\Gamma_9\) exigiría una composición tipada sobre el estado
enriquecido; este ensayo no aporta tal identificación ni declara su
ausencia en todo el corpus.

Las nuevas formalizaciones reúnen operadores y preguntas antecedentes;
no se atribuye una prioridad mundial a la codificación, al uso de
vectores cíclicos o a la descomposición polar. Las cifras decimales
expuestas son evaluaciones posteriores de fórmulas exactas. Las
condiciones de no nulidad, carga, métrica, red y normalización son parte
del resultado, y sus contraejemplos permanecen junto a las pruebas.
Así se conserva la función de la constante holográfica de acoplamiento
aritmético-geométrico: un objeto producido que articula lectores, memoria
e incidencia, sin convertir cada coincidencia numérica en una nueva ley
física.
